% BHU PYQ -- Manually transcribed & error-corrected
\documentclass[11pt,a4paper]{article}

\usepackage[a4paper,top=2.5cm,bottom=2.5cm,left=2.8cm,right=2.8cm,headheight=14pt]{geometry}
\usepackage{amsmath,amssymb}
\usepackage[T1]{fontenc}
\usepackage[utf8]{inputenc}
\usepackage{lmodern}
\usepackage{microtype}
\usepackage{fancyhdr}
\usepackage{enumitem}
\usepackage{hyperref}
\usepackage{lastpage}
\usepackage{parskip}

\hypersetup{colorlinks=true,urlcolor=black,linkcolor=black,
  pdftitle={CHB-602: Inorganic Chemistry-IV -- BHU PYQ},pdfauthor={Banaras Hindu University}}

\pagestyle{fancy}\fancyhf{}
\renewcommand{\headrulewidth}{0.4pt}\renewcommand{\footrulewidth}{0.4pt}
\fancyhead[L]{\small   \textit{Banaras Hindu University}}
\fancyhead[R]{\small   \textit{CHB-602: Inorganic Chemistry-IV}}
\fancyfoot[C]{\small Page  \thepage\ of \pageref{LastPage}}


\newlist{parts}{enumerate}{2}
\setlist[parts,1]{label=(\alph*),leftmargin=*,itemsep=5pt,topsep=3pt}
\setlist[parts,2]{label=(\roman*),leftmargin=*,itemsep=3pt,topsep=2pt}

\newcommand{\pts}[1]{\hfill   \textit{[#1]}}


\begin{document}


\begin{center}
  {\large\bfseries BANARAS HINDU UNIVERSITY}\\[2pt]
  {\normalsize B.Sc. (Hons.) Semester VI Examination, 2022-23}\\[6pt]
  \rule{0.8\linewidth}{0.5pt}\\[6pt]
  {\large\bfseries Paper No. CHB-602: Inorganic Chemistry-IV}\\[4pt]
  {\normalsize Time: 3 Hours\hfill Maximum Marks: 70}
\end{center}
\rule{\linewidth}{0.4pt}
\small



\noindent   \textit{Instructions: Answer five questions in total, including Question~1 which is compulsory.}

\normalsize
\bigskip




\noindent   \textbf{Question 1.} Answer all parts of the following: \pts{5 $ \times$ 2 = 10}

\begin{parts}
  \item Why are lower oxidation states more stable in $3 \text{d}$ metals and higher oxidation states are more stable in $4 \text{d}$ and $5 \text{d}$ metals?
  \item Which of the phenomena among paramagnetism, diamagnetism, and ferromagnetism does not occur in liquid or gaseous state and why?
  \item Explain the spin-orbit coupling conditions for $3 \text{d}$-metal ions.
  \item Differentiate between magnetically diluted and magnetically non-diluted substances.
  \item Using an Orgel diagram, assign electronic transitions to the absorption peaks for:
  
\begin{parts}
    \item $[ \text{CoCl}_4]^{2-}$ : $3300\, \text{cm}^{-1}$, $5800\, \text{cm}^{-1}$, $15000\, \text{cm}^{-1}$
    \item $[ \text{V(H}_2 \text{O)}_6]^{3+}$ : $17400\, \text{cm}^{-1}$, $25000\, \text{cm}^{-1}$, $37000\, \text{cm}^{-1}$
  \end{parts}
\end{parts}

\medskip\rule{\linewidth}{0.2pt}




\noindent   \textbf{Question 2.}

\begin{parts}
  \item Write the electronic configurations for the $5 \text{f}$ series of elements (actinides) and compare their oxidation states with those of the corresponding lanthanides. \pts{6}
  \item Discuss the separation of Lanthanides by the solvent extraction method. \pts{8}
\end{parts}

\medskip\rule{\linewidth}{0.2pt}




\noindent   \textbf{Question 3.}

\begin{parts}
  \item Explain the magnetic properties of tetrabutylammonium copper(II) chromium(III) oxalate. How does magnetic coupling occur? \pts{6}
  \item Calculate the spin-only and experimentally observed magnetic moments of the following ions: $ \text{Fe}^{3+}$, $ \text{Co}^{2+}$, $ \text{Ni}^{2+}$, and $ \text{Cu}^{2+}$. \pts{8}
\end{parts}

\medskip\rule{\linewidth}{0.2pt}




\noindent   \textbf{Question 4.}

\begin{parts}
  \item Compare complex formation tendencies of actinides with those of lanthanides. Give examples. \pts{6}
  \item Discuss the biological role of nitrogenase enzyme in nitrogen fixation. \pts{8}
\end{parts}

\medskip\rule{\linewidth}{0.2pt}




\noindent   \textbf{Question 5.}

\begin{parts}
  \item Describe the phenomenon of superexchange where individual magnetic dipoles of metal ions couple through the electrons of intervening ligands. \pts{8}
  \item How can $ \text{Ce}$ and $ \text{Eu}$ be separated from a mixture of other lanthanides using selective oxidation/reduction methods? \pts{6}
\end{parts}

\vfill
\rule{\linewidth}{0.4pt}

\begin{center}\small
  \href{https://bhu-exam.vercel.app/}{Click here to visit our website}\\[4pt]
  \href{https://me-aryan.vercel.app/}{Click here to visit our contributor}\\[4pt]
  \href{https://quashan-me.vercel.app/}{Click here to visit our contributor}
\end{center}

\end{document}
